\documentclass[10pt,a4paper]{amsart}
\usepackage[margin=19mm]{geometry}
\usepackage{amsmath,amssymb,mathtools}
\usepackage[scr=rsfs,cal=euler]{mathalfa}
\usepackage{xfrac}
\usepackage[unicode,hidelinks,bookmarks=false]{hyperref}
\newcommand{\ie}{i.e.\ }
\newcommand{\cf}{cf.\ }
\newcommand{\resp}{respectively\ }
\newcommand{\Subsection}{\subsection}
\providecommand{\nobreakdash}{\nobreak\hbox{-}}
\setlength{\emergencystretch}{2em}
\multlinegap0pt
\usepackage{bm}
\usepackage{amssymb}
\usepackage[all]{xy}
\newtheorem{theorem}{Theorem}
\title{{A class of bilateral weighted shift operators, and linear dynamics}}
\author{Bibhash Kumar Das, Aneesh Mundayadan}
\date{Source: arXiv:2602.22773v1 (CC BY 4.0)}
\begin{document}
\maketitle

\noindent\emph{Source and licence.} {A class of bilateral weighted shift operators, and linear dynamics}. Version arXiv:2602.22773v1 (CC BY 4.0), distributed by arXiv under the Creative Commons Attribution 4.0 International licence, http://creativecommons.org/licenses/by/4.0/, as recorded in the arXiv licence metadata for this exact version and shown on the abstract page https://arxiv.org/abs/2602.22773v1. Full licence notice: \texttt{licenses/arxiv-2602.22773-CC-BY-4.0.txt}.\par\smallskip

\noindent\textit{Editorial scope.} Counted unit: the article's theorem T (source0.tex lines 1374--1384). The article's complete original proof of it is delivered with it (source0.tex lines 1387--1453, one \texttt{proof} environment of 67 lines). No mathematics is written, repaired or re-derived: the statement and the proof are the article's own lines, in the article's own notation, and no retained line is substituted or re-flowed.  Retained uncounted as the source-local context the proof needs: lines 521--539; lines 778--795; lines 797--799; lines 822--843.  Nothing was omitted from any retained span, so the package carries no omission record.  No retained line contains a citation, so the wrapper carries no bibliography and no bibliography entry.  No retained line contains a figure environment, an image-inclusion command or drawing code, so no image is delivered and the package declares no asset dependency.  Editorial formatting only, in addition: an AMS \texttt{amsart} preamble that restates the article's own declarations and macros used by the retained lines, namely the environments \texttt{theorem} on separate counters, together with the standard packages those lines need.  The retained environments print in this document as Theorem 1 in order of appearance, whereas the article prints the same items with the numbering of its own full document.  The complete native source of the article as distributed in the arXiv e-print for this exact version is bundled unchanged as \texttt{sources/195350/source0.tex}.
   \begin{theorem}\label{decomposition1}
  Assume that $z^n\in \ell^p_{a,b}(\Omega_{r,R})$ for all $n\in \mathbb{Z}$, and 
  \begin{equation}\label{rr}
  \sup_{n\in \mathbb{Z}}~\left\lvert\frac{ w_{n+1}a_{n+1}}{a_{n}}\right\rvert<\infty.
  \end{equation}
  \begin{itemize}
  \item[(i)] If 
  \begin{equation}\label{ss}
  \sum_{i=1}^{+\infty}\sup_{n\in \mathbb{Z}}~~\left\lvert c_{n}\frac{b_{n}b_{n+1}\cdots b_{n+i-1}}{a_{n+1}a_{n+2}\cdots a_{n+i}}\right\rvert<\infty, ~\text{and}~ \sup_{n\in \mathbb{Z}}|c_n|<\infty,
  \end{equation}
  then $B_{w}$ is bounded on $\ell^{p}_{a,b}(\Omega_{r,R})$.
\item[(ii)] If 
\begin{equation}\label{tt}
\limsup_{ n\rightarrow +\infty} \left\lvert b_n/a_{n+1}\right\rvert<1,
\end{equation}
then the conditions in \textnormal{(i)} are satisfied and consequently, $B_w$ is bounded.
   \end{itemize}
   A similar result holds for $c_{0,a,b}(\Omega_{r,R}).$
   \end{theorem}

 \begin{equation}\label{matrix}
  [B_w^{\nu}]=
  \begin{bmatrix}
  \ddots&  \ddots & \ddots & \ddots & \ddots&\ddots&\ddots&\ddots&\ddots\\
   \ddots & A_{-1,\nu}&0&0&\cdots&0&0&0&\ddots\\
   \ddots & 0& A_{0,\nu}&0&\cdots&0&0&0&\ddots\\
  \ddots & 0& C_{0,\nu}&A_{1,\nu}&\cdots&0&0&0&\ddots\\
  \ddots & 0& 0&C_{1,\nu}& \cdots&0&0&0&\ddots\\
  \ddots & \vdots& \vdots&\vdots& \cdots&\vdots&\vdots&\vdots&\ddots\\
\ddots& 0 & 0&0 & \cdots&A_{\nu-2,\nu}&0&0 &\ddots\\
\ddots& 0 & 0 & 0&\cdots&C_{\nu-2,\nu}&A_{\nu-1,\nu}&0&\ddots\\
  \ddots &0&\boxed{\boldsymbol{0}} &0&\cdots &0&C_{\nu-1,\nu}& A_{\nu,\nu}&\ddots\\
	\ddots &0&0& 0&\cdots&0&-C_{\nu-1,\nu}\frac{b_{0}}{a_{1}}& C_{\nu,\nu}&\ddots\\
	\ddots  & 0&0&0 &\cdots&0&C_{\nu-1,\nu}\frac{b_{0}b_{1}}{a_{1}a_{2}}&-C_{\nu,\nu}\frac{b_{1}}{a_{2}}&\ddots\\
 \ddots  & \ddots&0&0&\cdots&0&-C_{\nu-1,\nu}\frac{b_{0}b_{1}b_{2}}{a_{1}a_{2}a_{3}}&C_{\nu,\nu}\frac{b_{1}b_{2}}{a_{2}a_{3}}&\ddots\\
	  \ddots & \ddots&\ddots&\ddots &\ddots&\ddots&\ddots&\ddots&\ddots 
	\end{bmatrix}.
 \end{equation}

\begin{equation}\label{cc}
   [B_{w}^{\nu}]u=\sum_{j=-\infty}^{+\infty} \alpha_{j,\nu} ~ e_{j},
\end{equation}

 \begin{theorem}\label{hyper5}
      The following hold for the bilateral weighted backward shift $B_{w}$ acting on  $\ell^{p}_{a,b}(\Omega_{r,R})$, $1\leq p<\infty,$ assuming that the conditions \eqref{rr} and \eqref{tt} of Theorem \ref{decomposition1} are satisfied:
     \begin{enumerate}
      \item[(i)] $B_{w}$ is hypercyclic on $\ell^{p}_{a,b}(\Omega_{r,R})$ if and only if
         \begin{equation}\label{hyperc}
              \limsup_{\nu\rightarrow +\infty}\left\lvert w_{n+1}\cdots w_{n+\nu}{a_{n+\nu}}   
 \right\rvert=\infty= \limsup_{\nu\rightarrow +\infty} \left\lvert \frac{a_{n-\nu}}{w_{n}\cdots w_{n-\nu+1}} \right\rvert,~~ \forall~~ n\in \mathbb{N}.
        \end{equation}
         \item[(ii)] $B_{w}$ is topologically mixing on $\ell^{p}_{a,b}(\Omega_{r,R})$ if and only if
         \begin{equation}\label{mixing}
         \lim_{\nu\rightarrow +\infty}\left\lvert w_{n+1}\cdots w_{n+\nu}{a_{n+\nu}}   
 \right\rvert=\infty= \lim_{\nu\rightarrow +\infty} \left\lvert \frac{a_{n-\nu}}{w_{n}\cdots w_{n-\nu+1}}\right\rvert,~~ \forall~~ n\in \mathbb{N}.
         \end{equation}

         \item[(iii)] $B_{w}$ is supercyclic on $\ell^{p}_{a,b}(\Omega_{r,R})$ if and only if
         \begin{equation}\label{superc}
         \limsup_{\nu\rightarrow +\infty} \left\lvert 
         \frac{w_{n+1}\cdots w_{n+\nu}}{w_{n}\cdots w_{n-\nu+1}}
          a_{n+\nu}a_{n-\nu}\right\rvert=\infty, \hskip .5cm \forall~~ n\in \mathbb{N}.
         \end{equation}
        \end{enumerate}
    \end{theorem}

\begin{theorem}\label{T}
    Assume that, for $1\leq p<\infty,$ 
    \begin{equation}\label{yy}
\sup_{n\in \mathbb{Z}}~\left\lvert\frac{ w_{n+1}a_{n+1}}{a_{n}}\right\rvert<\infty, \hspace{.4cm} \limsup_{n\rightarrow +\infty}\left \lvert \frac{b_{n}}{a_{n+1}}\right \rvert<1,~~~~ \text{and} \hspace{.4cm} \sup_{n\geq 1}~\left\lvert\frac{ w_{0}\cdots w_{-n+1}}{a_{-n}}\right\rvert<\infty.
\end{equation}
Then, the following statements are equivalent:
      \begin{itemize}
          \item[(i)] $B_{w}$ is not hypercyclic on $\ell^p_{a,b}(\Omega_{r,R}).$
             \item[(ii)] For each $u \in \ell^p(\mathbb{Z}),$ one has that $[B_{w}^{\nu}]u\to 0$ in $\mathbb{C}^{\mathbb{Z}}$ as $\nu\to +\infty$.
      \end{itemize}
\end{theorem}

\begin{proof}
 For $\nu\geq 1$ and $u=(\lambda_{n})_{n\in \mathbb{Z}}\in \ell^{p}(\mathbb{Z})$, recalling the matrix of $B_{w}^{\nu}$ from (\ref{matrix}) and the action of $[B_{w}^{\nu}]u$ from \eqref{cc}, we have
\begin{equation}\label{ccc}
   [B_{w}^{\nu}]u=\sum_{n=-\infty}^{+\infty} \alpha_{n,\nu} ~ e_{n},
\end{equation}
where $\{e_{n}\}_{n\in \mathbb{Z}}$ is the standard basis in $\ell^p(\mathbb{Z})$, and 
$\alpha_{n,\nu}$ is as given in \eqref{cc}.
Now, assume that $B_{w}$ is not hypercyclic on $\ell^p_{a,b}(\Omega_{r,R})$. By Theorem \ref{hyper5}, we have 
\begin{equation}\label{ddd}
M:=\sup_{\nu\geq 1}\left|w_{n+1}\cdots w_{n+\nu}{a_{n+\nu}}\right|<\infty,\quad\forall n\in \mathbb{N}.
\end{equation}
We will show that 
\[
\lim_{\nu\to+\infty}\alpha_{n,\nu}=0 , \qquad n\in\mathbb{Z},
\]
by examining the individual terms appearing in $\alpha_{n,\nu}$. First, note that the hypothesis $\limsup_{\nu\to+\infty}|b_\nu/a_{\nu+1}|<1$ yields some numbers $r<1$ and $N\geq 1$ such that  
\begin{equation}\label{dd}
\big|b_\nu/a_{\nu+1}\big|<r, \qquad \nu\geq N.
\end{equation}

\noindent \textbf{Case $n\le -\nu$:} Since $A_{\nu+n,\nu}
=w_{\nu+n}\cdots w_{n+1}\frac{a_{\nu+n}}{a_{n}}$, the assumptions \eqref{yy} give that the family $\{A_{\nu+n,\nu}:\nu\ge 1\}$ is bounded for each fixed $n\le -\nu$.

\noindent{\bf Case $-\nu+1\le n\le 0$:}
We have
$
C_{\nu+n-1,\nu}
   = w_{\nu+n}\cdots w_{n+1} \frac{b_{\nu+n-1}}{a_{n}},
~\text{and}~
A_{\nu+n,\nu}=
   w_{\nu+n}\cdots w_{n+1}\frac{a_{\nu+n}}{a_{n}}.
$
Using the equations \eqref{yy}, \eqref{ddd}, and the estimate 
$|b_\nu/a_{\nu+1}|<r$, both families
$
\{C_{\nu+n-1,\nu}:\nu\ge 1\}$ and 
$\{A_{\nu+n,\nu}:\nu\ge 1\}
$
become bounded for fixed $n$.

\noindent{\bf Case $n\ge 1$:}
 Note that $A_{\nu+n,\nu}=w_{\nu+n}\cdots w_{n+1}\frac{a_{\nu+n}}{a_{n}}$,
and so, $\{A_{\nu+n,\nu}:\nu\geq 1\}$ is bounded for fixed $n$. Finally, we have
\[
\begin{aligned}
|C_{\nu+n-1,\nu}|
&\le
\left|\frac{w_{\nu+n}\cdots w_{n+1}a_{\nu+n}}{a_{n}}\right|
   \left|\frac{b_{\nu+n-1}}{a_{\nu+n}}\right|+
\left|\frac{w_{\nu+n-1}\cdots w_{n}a_{\nu+n-1}}{a_{n-1}}\right|
   \left|\frac{b_{n-1}}{a_{n}}\right|.
\end{aligned}
\]
Both the terms are bounded, and therefore  
$
\{C_{\nu+n-1,\nu}:\nu\geq 1\}
$
is bounded. Using \eqref{dd} and H\"{o}lder's inequality, in all three cases we find that, for large $\nu$
$$
|\alpha_{n,\nu}|\leq\begin{cases}
   C_{1}|\lambda_{\nu+n}|,\hspace{3.4cm} \text{if } n\leq -\nu,\\ 
   C_{2}(|\lambda_{\nu+n-1}|^{  p}+|\lambda_{\nu+n}|^{p})^{\frac{1}{p}},\hspace{.6cm} \text{if }-\nu+1\leq n\leq 0,\\ 
    C_{3} \left(\sum_{j\geq \nu-1}^{\nu+n}|\lambda_j|^p\right)^{1/p}, \hspace{1.2cm} \text{if }n\geq 1,\\ 
    \end{cases}
    $$
where $C_{1},C_{2},C_{3}$ are constants depending only on $r$. Combining all these inequalities, we conclude that the coefficients in \eqref{ccc} converge to $0$, as $\nu\rightarrow +\infty$.
\end{proof} 

\end{document}
